\section{Related Work}
ReAct provides an example of a language model interleaving reasoning with actions that retrieve external information \cite{yao2023react}. In the proposed system, the agent selects from predefined tools, and the application validates the requested parameters and retrieved records. The UCI Occupancy Detection dataset provides indoor measurements with occupancy labels \cite{uci_occupancy}. In this project, it is considered for checking the preprocessing workflow. Labels for appliance operation will be collected separately through operation logs and user feedback. The CBE Thermal Comfort Tool illustrates the use of multiple inputs in thermal comfort assessment \cite{tartarini2020cbe}. Our single-room prototype will report the limits of recommendations based on temperature and relative humidity measurements.


Prior smart-home research, including CASAS, provides a precedent for collecting residential sensor observations and analyzing daily activities \cite{cook2013casas}. Our target labels describe actions such as operating an appliance or opening a window, so the evaluation requires records collected in the monitored room. Statistical changepoint methods locate changes in a time series \cite{killick2012changepoints}. The proposed prototype uses a sustained-change rule to identify candidate events and links those detections to separately collected operation reports.

Ecological momentary assessment motivates collecting reports close to an event in everyday settings \cite{shiffman2008ema}. Event-triggered prompts ask about possible actions while the context is recent. Because replies can be delayed or uncertain, the system records event onset, prompt time, reply time, and reported operation time separately. Table~\ref{tab:literature} compares the related studies. MQTT and PROV-O provide technical specifications for sensor messaging and provenance representation, respectively \cite{oasis2019mqtt,w3c2013provo}.

\begin{table*}[t]
\caption{Article Summary Table}
\label{tab:literature}
\centering\footnotesize
\setlength{\tabcolsep}{3pt}
\begin{tabular}{@{}p{.10\textwidth}p{.11\textwidth}p{.12\textwidth}p{.135\textwidth}p{.135\textwidth}p{.14\textwidth}p{.14\textwidth}@{}}
\hline
Study & Objective & Method & Data or setting & Main finding or contribution & Limit for this study & Use in this proposal \\
\hline
ReAct \cite{yao2023react} &
Integrate reasoning and action &
Alternate reasoning with tool calls &
Question answering, fact verification, and interactive tasks &
Actions retrieve evidence that can inform subsequent reasoning &
Benchmarks differ from room sensing and weather advice &
Route queries to predefined tools and validate evidence \\
UCI Occupancy Detection \cite{uci_occupancy} &
Identify room occupancy &
Binary classification with photo-derived labels &
Office temperature, humidity, light, and carbon dioxide &
Provides time-series training and test sets with occupancy labels &
Occupancy labels do not distinguish appliance or window actions &
Check preprocessing and collect separate operation logs \\
CBE Thermal Comfort Tool \cite{tartarini2020cbe} &
Assess thermal comfort &
Web implementation of comfort models &
Environmental and personal inputs &
Makes model-based comfort calculations available online &
This prototype observes only temperature and humidity &
Keep comfort advice conditional on available evidence \\
CASAS \cite{cook2013casas} &
Support smart-home activity studies &
Deploy residential sensing and activity analysis &
Smart-home sensor observations &
Describes a platform for observing everyday activities &
Activity categories do not supply labels for this room &
Collect and label actions in the monitored room \\
Killick et al.\ \cite{killick2012changepoints} &
Locate multiple changes &
Pruned exact minimization of a penalized cost &
Simulated time series and application data &
Linear computational cost under stated conditions &
A detected change does not identify its physical cause &
Use sustained change to propose candidate events \\
Shiffman et al.\ \cite{shiffman2008ema} &
Study experience in everyday settings &
Review of repeated reports close to events &
Ambulatory behavioral research &
Describes methods for reducing retrospective recall demands &
Completeness depends on whether participants respond &
Record prompt, reply, and reported action times separately \\
Roberts et al.\ \cite{roberts2017cv} &
Validate models with dependent observations &
Review and comparison of blocked validation &
Examples with structured ecological data &
Blocking can reduce bias from dependence across splits &
The split must match the intended prediction setting &
Keep events intact and hold out later events \\
\hline
\end{tabular}
\end{table*}


